% ECONOMICS_PROBLEM: {"id": "I18-365", "title": "Optimal policy toward addiction and opioids", "jel_code": "I18", "source_id": "OP-0057"}
% !TeX program = xelatex
\documentclass[11pt,a4paper]{article}
\usepackage[margin=23mm]{geometry}
\usepackage{fontspec}
\setmainfont{Latin Modern Roman}
\usepackage{amsmath,amssymb,mathtools}
\usepackage{xcolor,enumitem,booktabs,longtable,array,xurl,hyperref}
\definecolor{muted}{HTML}{53636C}
\hypersetup{colorlinks=true,urlcolor=blue,linkcolor=blue}
\setlength{\parindent}{0pt}
\setlength{\parskip}{5pt}
\newcommand{\R}{\mathbb R}
\newcommand{\E}{\mathbb E}
\newcommand{\N}{\mathbb N}
\newcommand{\ind}{\mathbf 1}
\newcommand{\Var}{\operatorname{Var}}
\newcommand{\Cov}{\operatorname{Cov}}
\newcommand{\supp}{\operatorname{supp}}
\DeclareMathOperator*{\argmin}{arg\,min}
\DeclareMathOperator*{\argmax}{arg\,max}
\newcommand{\entrymeta}[3]{{\sffamily\small\color{muted}\textbf{\texttt{#1}}\quad|\quad #2\par #3\par}}
\newcommand{\Problemref}[1]{\texttt{#1}}
\newcommand{\JELref}[2]{\texttt{#2} (JEL \texttt{#1})}
\begin{document}
\section*{Optimal policy toward addiction and opioids}
\textbf{Problem ID:} \texttt{I18-365}\quad \textbf{JEL:} \texttt{I18}\par
% BEGIN ECONOMICS STATEMENT
\label{problem:OP-0057}
\label{atlas:prob:H7}
\noindent\textbf{Original atlas identifier:} H7; entry 57. \textbf{Field:} Health economics. \textbf{Original tractability label:} Medium.

\noindent\textbf{Classification:} Parameterized theoretical/empirical research agenda; not a certified unsolved theorem.

\subsection*{Definitions and assumptions}
Let state $s_t$ contain addiction stock, health, pain severity, treatment status, and legal/illicit market conditions. Individual choices $x_t=(x_t^{L},x_t^{I},r_t)$ denote legal use, illicit use, and treatment participation. A model $m$ specifies state transition $P_m(s_{t+1}\mid s_t,x_t,p_t)$, overdose risk, prescribing/provider behavior, illicit supply, and equilibrium selection. A policy $p_t$ comprises legal prices/taxes, prescription rules, treatment access, and enforcement. Fix discount factor $\beta\in(0,1)$ and welfare criterion $u_i^{W}(s_t,x_t)$, including pain relief and health harms. This criterion may differ from the decision utility generating behavior; if present bias is modeled, define it separately. Let $C_t(p;m)$ be real treatment, administration, and enforcement costs and $E_t$ external harms, all monetized consistently.

\subsection*{Formal research question}
\emph{Editorial formulation: the mathematical target below is not a theorem quoted from the references.}

Among legally feasible feedback rules $p_t=\pi(s_t^{\rm obs})$ depending only on observed admissible information, characterize
\[
 \sup_{\pi\in\Pi}\mathbb E_m\sum_{t=0}^{\infty}\beta^t
 \left[\sum_iw_i u_i^{W}(s_t,x_t)-C_t(\pi;m)-E_t(\pi;m)\right],
\]
with $w_i\geq0$ and a stated public financing constraint. Use a common monetary normalization for welfare and costs; taxes are transfers within the specified social boundary. Determine which supply, prescribing, and treatment shocks identify substitution between legal and illicit use, addiction persistence, and mortality. Compare joint policy packages over an identified model class, reporting welfare bounds and $\varepsilon$-optimal rules rather than extrapolating a single-market demand elasticity.

\subsection*{Interpretation and scope}
Restriction-induced reductions in legal prescriptions do not alone establish fewer deaths: users can switch substance, potency, supplier, or treatment. Changes in pain relief and untreated pain also belong in welfare. Rational-addiction and present-biased models imply different internality corrections, so the planner's normative benchmark must be stated rather than chosen implicitly from one model's demand equation. Illicit-market entry and quality are endogenous, and enforcement may displace geography rather than eliminate supply. Mortality records, prescriptions, and treatment records need linkage and uncertainty about unobserved illicit exposure. A policy shock's exclusion restriction must rule out simultaneous changes in enforcement or treatment if a single mechanism is to be identified. The extension studies joint dynamic policy with clinical and market substitution, beyond existing addiction and opioid-reformulation results.

\subsection*{Source audit and status}
The three original references are identifiable; the atlas's author spelling is corrected to Botond Köszegi. [S1] and [S2] are distinct behavioral benchmarks, and [S3] provides substitution evidence in a particular reformulation setting. They do not establish that a single addiction model is universally valid or certify the proposed joint-policy problem as open in 2026.

\subsection*{References}
\begin{enumerate}[label={[S\arabic*]},leftmargin=*,itemsep=0.6em]
\item Gary S. Becker and Kevin M. Murphy (1988). A Theory of Rational Addiction. Journal of Political Economy 96(4), 675--700. DOI: 10.1086/261558.
\url{https://www.journals.uchicago.edu/doi/10.1086/261558}
\newline\emph{Locator and support limits:} Abstract and model: dynamically consistent addiction benchmark; not evidence that all addiction is rational.
\item Jonathan Gruber and Botond Köszegi (2001). Is Addiction “Rational”? Theory and Evidence. The Quarterly Journal of Economics 116(4), 1261--1303. DOI: 10.1162/003355301753265570.
\url{https://academic.oup.com/qje/article-abstract/116/4/1261/1903215}
\newline\emph{Locator and support limits:} Abstract and model: time inconsistency and addiction. Corrected spelling to Köszegi.
\item Abby Alpert, David Powell, and Rosalie Liccardo Pacula (2018). Supply-Side Drug Policy in the Presence of Substitutes: Evidence from the Introduction of Abuse-Deterrent Opioids. American Economic Journal: Economic Policy 10(4), 1--35. DOI: 10.1257/pol.20170082.
\url{https://www.nber.org/papers/w23031}
\newline\emph{Locator and support limits:} Abstract and article: abuse-deterrent opioids and substitution; treatment access and joint optimal policy are extensions.
\end{enumerate}

\subsection*{Common conventions: Source mathematical conventions}

Unless an entry states otherwise, agent and firm sets are finite. State and action spaces are measurable, strategies and outcome maps are measurable, probability laws are specified, and expectations exist for the targets under discussion. Infinite-horizon discount factors lie in $(0,1)$ and discounted objectives are finite. Norms, tolerances and complexity input sizes must be specified for computational comparisons. Constants or primitives that appear locally are inputs, not empirical facts asserted by this catalogue.

\textbf{Identification.} A statistical model $m\in\mathcal M$ implies an observable law $P_m$ and target $\theta(m)$. At law $P$, its identified set is
\[
\Theta(P;\mathcal M)=\{\theta(m):m\in\mathcal M,\ P_m=P\}.
\]
Point identification holds when this set is a singleton, or equivalently when $P_{m_1}=P_{m_2}$ implies $\theta(m_1)=\theta(m_2)$ for all compatible models. Sharp bounds mean bounds attained, or approached when the boundary is not attained, by compatible models. Writing this set is a definition, not a solution: the substantive task is to establish informative restrictions and derive or estimate the set. An unrestricted-confounding model may give only trivial bounds.

\textbf{Inference.} Identification, estimability and valid uncertainty quantification are separate questions. Fix $\alpha\in(0,1)$. A confidence procedure $C_n$ over a declared law class $\mathcal P$ has asymptotic uniform coverage at level $1-\alpha$ when
\[
\liminf_{n\to\infty}\inf_{P\in\mathcal P}\Pr_P\{\theta(P)\in C_n\}\geq1-\alpha.
\]
For set-identified targets the coverage event must be defined separately, for example coverage of the full identified set. If an entry uses identification-indexed models instead of uniquely indexed laws, the same coverage convention is applied over those models. Pointwise limits do not establish uniform validity, and asymptotic validity does not guarantee finite-sample precision. Sample sizes, dependence structures and weak-identification sequences are part of each application.

\textbf{Counterfactuals.} $Y_i(a)$ is a potential outcome under a specified intervention, including its timing, implementation and versions. It is not $Y_i$ conditional on voluntarily choosing $a$. A full intervention vector permits interference; shorthand scalar treatment notation does not silently impose no interference. Assignment, selection, attrition, anticipation and measurement must be specified. A claim about an instrument's local effect is not automatically a population policy effect.

\textbf{Economic equilibrium.} A counterfactual model includes best responses, constraints, feasibility, market clearing, state transitions and expectations consistent with the stated information. When several equilibria exist, either a declared selector is an input or the target is set-valued. Comparisons across policy regimes must use the stated selection convention. Reduced-form relationships are not presumed invariant after an equilibrium-changing intervention.

\textbf{Optimization and welfare.} Optimization is over the stated feasible set. If attainment is not established, $\sup$ or $\inf$ and $\varepsilon$-optimal choices replace an asserted maximizer or minimizer. For example, with value $V=\sup_{a\in\mathcal A}W(a)<\infty$, an $\varepsilon$-optimal set is $\{a\in\mathcal A:W(a)\geq V-\varepsilon\}$ for $\varepsilon>0$. Benefits, costs, risk and health or environmental outcomes can be aggregated only using declared units and conversion weights. Interpersonal utility weights, the treatment of fiscal transfers and foreign profits, discounting, and the behavioral welfare criterion are explicit inputs. A comparative-static derivative additionally requires existence and appropriate differentiability of the selected counterfactual.

\textbf{Sources.} Local labels [S1], [S2], and so forth restart in every question. Locators identify the relevant abstract, model, section or result at the level actually checked; a bibliographic or abstract-level verification is not represented as inspection of an entire proof. Working papers are distinguished from later publications and changing versions. Source summaries are paraphrased. All URL checks and editorial formulations refer to the audit date stated above.
% END ECONOMICS STATEMENT
\end{document}
